\subsection{可忽略集}

定义 2. --- 给定一个测度 \(\mu\) 于局部紧空间 \(X\) 上，子集 \(A\) ⊂ \(X\) 称为关于测度 \(\mu\) 可忽略的，如果 \(|\mu|^*(A) = 0\) 。

我们也说 A 是 \(\mu\) -可忽略的，或在不致混淆时简称可忽略的。说特征函数 \(\varphi_{A}\) 可忽略是一回事。

命题 4. --- 可忽略集的每个子集都是可忽略的；可忽略集的可数并是可忽略的。

这是命题 1 与命题 2 的直接推论。

例子. --- 设 \(\mu\) 是 R 上的 Lebesgue 测度。每个化为一点的集合 \(\{x_{0}\}\) 都是可忽略的（参见 §1, No.~3, 注记 1）。由此推出 R 的每个可数子集关于 Lebesgue 测度都是可忽略的。该命题的逆命题不成立（§4, 习题 4 b)）。

命题 5. --- \(\mu\) 的支集 S 的余集是 X 中最大的可忽略开集。

事实上，由命题 3，要使开集 G 可忽略，必须且只须 \(G \cap S = \varnothing\) ，即 \(G \subset CS\) 。

